\section{Introduction}
\label{sec:01_introdoction}

Fine-grained visual categorization (FGVC) aims to distinguish highly similar subordinate categories where inter-class variance is frequently dominated by large changes in pose, illumination, background clutter, and scale. Pet-breed recognition represents a classic FGVC challenge: distinct breeds often differ only in localized facial geometry, coat texture, or anatomical proportions. In this work, we address a constrained 20-class pet-breed recognition task subject to specific operational rules. The target output space comprises 10 cat breeds (classes $0$ to $9$) and 10 dog breeds (classes $10$ to $19$), supplemented by a mandatory reject class ($-1$) for images containing no valid target species or unresolvable confounders. 

Furthermore, two explicit constraints dictate the task:
\begin{enumerate}
    \item \textbf{Largest Target Rule:} When multiple target animals co-occur within a single frame, the ground-truth label strictly corresponds to the spatially largest animal by bounding-box area.
    \item \textbf{Runtime Budget:} To ensure real-time viability, the full end-to-end evaluation pipeline must execute with a mean inference latency strictly below $5.0$ seconds per image.
\end{enumerate}

Monolithic, single-stage classification networks frequently suffer in unconstrained settings because visual representations can easily exploit background shortcuts or misallocate capacity to secondary non-target animals. To overcome these limitations, we propose a two-stage cascade architecture. A COCO-pretrained \texttt{YOLOv6} \cite{li2022yolov6} detector first localizes candidates, resolves spatial competition by selecting the largest target bounding box, and zero-masks secondary animal detections to prevent context leakage. The primary crop is resized to $224\times224$ pixels and routed by predicted species (cat vs. dog) to one of two dedicated \texttt{EfficientNetV2-S} \cite{tan2021efficientnetv2} classification backbones. This divide-and-conquer strategy reduces the problem space to two simpler $10$-class sub-problems, eliminates inter-species capacity competition, and delegates image rejection entirely to the detector thresholding stage.

The primary contributions of this work are as follows:
\begin{itemize}
    \item \textbf{End-to-End Cascade Pipeline:} An inference pipeline combining \texttt{YOLOv6} and zero-out spatial masking with species-routed \texttt{EfficientNetV2-S} classifiers, fully satisfying the largest-target rule and rejection constraints.
    \item \textbf{Curated Multisource Dataset:} An integrated dataset of $35,423$ verified cat and dog images, constructed by merging Oxford-IIIT Pet \cite{parkhi2012catsdogs}, Stanford Dogs \cite{khosla2011novel}, The-Cat-API, Kaggle Dog, and targeted web scrapes to mitigate severe class imbalances.
    \item \textbf{Regularization \& Convergence Dynamics:} An empirical study showing that spatial patch mixing (\texttt{CropMix}) acts as a decisive regularizer when training from scratch, driving dog breed accuracy from $78.66\%$ to $85.02\%$. Crucially, we demonstrate that extended training epochs yield no further capacity gains and exacerbate overfitting, validating the necessity of threshold-driven early stopping.
    \item \textbf{Explainability \& Latency Benchmarking:} Grad-CAM diagnostic maps confirming that spatial masking suppresses background shortcut learning, alongside profiling proving a mean per-image latency of $200.63~\text{ms}$ ($4.01\%$ of the allowed time budget).
\end{itemize}

All source code, trained checkpoints, and reproducibility scripts are publicly available at: \href{https://github.com/NordlichtS/CVDL26-drop-table-catdogs}{CVDL26-drop-table-catdogs}.

The remainder of this paper is structured as follows: Section~\ref{sec:related_work} reviews foundational literature on object detection, fine-grained classification, and explainability. Section~\ref{sec:method} presents our two-stage cascade architecture, zero-out masking strategy, and routing interface. Section~\ref{sec:experiments} details the multi-source dataset integration, training dynamics, spatial regularization experiments, Grad-CAM interpretability, and runtime benchmarks. Finally, Section~\ref{sec:conclusion} summarizes our core findings and outlines future directions.